\documentclass[11pt]{article}
\usepackage[a4paper,margin=1in]{geometry}
\usepackage{kotex}
\usepackage{amsmath}
\usepackage{amssymb}
\usepackage{booktabs}
\usepackage{array}
\usepackage{xcolor}
\usepackage{listings}

\lstset{
  basicstyle=\ttfamily\small,
  breaklines=true,
  columns=fullflexible,
  frame=single,
  numbers=left,
  numberstyle=\tiny,
  keywordstyle=\color{blue!60!black},
  commentstyle=\color{green!40!black}
}

\begin{document}
\sloppy

\section*{PSI-Innerproduct의 RsCpsi 및 B2A-OT 구현 정리}

본 문서는 새로 추가한 \texttt{PSI\_Innerproduct} 구현을 설명한다.
현재 구현은 RsCpsi를 이용해 sender의 associated data를 교집합 위치 기준으로
정렬하고, B2A-OT를 통해 Boolean share가 곱해진 \(\mathbb{Z}_p\) arithmetic
share를 생성한다. 최종 output은 각 bin \(i\)에 대해 다음을 만족한다.

\[
  a_i^S + a_i^R = b_i \cdot D_i^Y \pmod p.
\]

여기서 \(b_i\)는 해당 위치가 교집합에 속하는지를 나타내는 indicator이고,
\(D_i^Y\)는 sender가 가진 associated data이다. 즉 교집합 위치에서는
associated data share가 유지되고, 교집합이 아닌 위치에서는 결과가 0이 되도록
만든다.

\section*{추가 파일과 역할}

\begin{table}[h]
\centering
\begin{tabular}{ll}
\toprule
파일 & 역할 \\
\midrule
\texttt{volePSI/PSI\_Innerproduct.h} & wrapper class와 B2A helper API 선언 \\
\texttt{volePSI/PSI\_Innerproduct.cpp} & IKNP OT 기반 B2A-OT 및 \(\mathbb{Z}_p\) 연산 구현 \\
\texttt{volePSI/CMakeLists.txt} & \texttt{PSI\_Innerproduct.cpp} 빌드 대상 추가 \\
\bottomrule
\end{tabular}
\end{table}

\section*{전체 구조}

Sender는 식별자 집합 \(Y\)와 associated data \(D^Y\)를 입력하고, receiver는
식별자 집합 \(X\)를 입력한다. 두 참여자는 먼저 prime mode의 RsCpsi를 수행한다.
RsCpsi는 sender의 associated data를 receiver의 cuckoo bin 기준으로 정렬된 share로
만든다. 이 단계 이후 각 bin \(i\)에 대해 양쪽은 다음 값을 가진다.

\[
  b_i = b_i^S \oplus b_i^R.
\]

\(b_i=1\)이면 해당 bin은 교집합 원소에 대응하고, \(b_i=0\)이면 교집합이 아닌
위치이다. 또한 prime mode RsCpsi의 value output은 sender associated data의
additive share로 해석된다. 교집합 위치에서는 다음 관계가 성립한다.

\[
  D_i^Y = w_i^S + w_i^R \pmod p.
\]

따라서 필요한 값은 단순히 \(D_i^Y\)의 share가 아니라 indicator가 곱해진 값이다.

\[
  b_i \cdot D_i^Y = b_i \cdot (w_i^S+w_i^R) \pmod p.
\]

이를 위해 구현은 B2A-OT를 두 번 수행한다. 첫 번째 B2A-OT는 sender가 가진
\(w_i^S\)에 대해 \(b_i w_i^S\)의 arithmetic share를 만든다. 두 번째 B2A-OT는
receiver가 가진 \(w_i^R\)에 대해 \(b_i w_i^R\)의 arithmetic share를 만든다.
마지막으로 두 결과를 \(\mathbb{Z}_p\)에서 더해 \(b_iD_i^Y\)의 share를 만든다.

\section*{Config 구조}

\texttt{PsiInnerproductConfig}는 field와 길이 정보를 관리한다. 기본 prime은
RsCpsi prime mode와 동일한 \(4294967291\)이다.

\begin{lstlisting}[language=C++]
struct PsiInnerproductConfig
{
    u64 mPrime = RsCpsiDefaultPrime;
    u64 mStatSecParam = 40;
    u64 mNumThreads = 1;

    u64 dataByteLength() const { return RsCpsiDataByteLength(mPrime); }
    u64 shareByteLength() const { return RsCpsiPrimeByteLength(mPrime); }
};
\end{lstlisting}

기본 설정에서는 input associated data가 16 bit이고, arithmetic share는 32 bit이다.
따라서 sender input associated data는 2 byte로 저장되고, RsCpsi 및 B2A 이후의
share는 4 byte로 저장된다. 내부 계산에서는 \texttt{u64}를 사용하지만 matrix에는
필요한 byte length만 기록한다.

\begin{table}[h]
\centering
\begin{tabular}{lll}
\toprule
항목 & 기본값 & 의미 \\
\midrule
\texttt{mPrime} & \(4294967291\) & 연산 field \(\mathbb{Z}_p\) \\
\texttt{dataByteLength()} & 2 byte & sender associated data 입력 길이 \\
\texttt{shareByteLength()} & 4 byte & CPSI/B2A arithmetic share 길이 \\
\texttt{mStatSecParam} & 40 & RsCpsi security parameter \\
\texttt{mNumThreads} & 1 & protocol thread 수 \\
\bottomrule
\end{tabular}
\end{table}

\section*{RsCpsi Wrapper 구현}

\texttt{PsiInnerproductSender}와 \texttt{PsiInnerproductReceiver}는 내부적으로
각각 \texttt{RsCpsiSender}, \texttt{RsCpsiReceiver}를 가진다. 초기화할 때
\texttt{ValueShareType::prime}을 사용하도록 고정하였다. 따라서 wrapper는 XOR mode나
add32 mode가 아니라, inner product 계산에 필요한 \(\mathbb{Z}_p\) share를 얻는
경로만 사용한다.

Sender 쪽 input은 \texttt{span<block> identifiers}와 \texttt{span<const u64>
associatedData}이다. associated data는 \texttt{dataByteLength()}만큼 matrix에
기록한 뒤 RsCpsi에 전달한다.

\begin{lstlisting}[language=C++]
oc::Matrix<u8> values(associatedData.size(), mConfig.dataByteLength());
setPsiIpPrimeValues(values, associatedData, mConfig.dataByteLength());
co_await mCpsi.send(identifiers, values, share, chl);
\end{lstlisting}

Receiver 쪽은 associated data 없이 identifier만 입력한다. RsCpsi receiver output에는
receiver의 bin mapping, Boolean flag share, 그리고 sender value share의 receiver 몫이
들어 있다.

\begin{lstlisting}[language=C++]
co_await mCpsi.receive(identifiers, share, chl);
\end{lstlisting}

\section*{B2A-OT 목표}

B2A-OT의 입력은 한쪽이 가진 value \(x\in\mathbb{Z}_p\)와 양쪽이 가진 Boolean
share \(b^{owner}, b^{choice}\)이다. 두 share의 XOR가 실제 bit이다.

\[
  b = b^{owner} \oplus b^{choice}.
\]

목표는 두 참여자가 \(z^{owner}, z^{choice}\in\mathbb{Z}_p\)를 얻고 다음을 만족하게
하는 것이다.

\[
  z^{owner} + z^{choice} = b\cdot x \pmod p.
\]

현재 구현에서는 libOTe의 IKNP 1-out-of-2 OT extension을 사용한다. value \(x\)를
가진 쪽이 OT sender가 되고, 자기 Boolean share를 선택 bit로 사용하는 쪽이 OT
receiver가 된다.

\begin{table}[h]
\centering
\begin{tabular}{ll}
\toprule
함수 & 역할 \\
\midrule
\texttt{psiIpB2aValueOwner} & value \(x\)를 가진 쪽, OT sender 역할 \\
\texttt{psiIpB2aChoiceOwner} & 선택 bit를 가진 쪽, OT receiver 역할 \\
\texttt{psiIpAddPrimeShares} & 두 B2A 결과를 \(\mathbb{Z}_p\)에서 더함 \\
\bottomrule
\end{tabular}
\end{table}

\section*{B2A-OT 메시지 구성}

Value owner는 각 row와 각 field element마다 랜덤 \(r\leftarrow\mathbb{Z}_p\)를
샘플링한다. 그리고 자기 arithmetic share를 \(r\)로 둔다. 그러면 choice owner가
받아야 하는 값은 실제 bit \(b\)에 따라 다음과 같다.

\[
  b=0: -r, \qquad b=1: x-r.
\]

하지만 value owner는 choice owner의 Boolean share를 모르므로 실제 \(b\)를 직접 알 수
없다. 그래서 value owner는 자기 Boolean share \(b^{owner}\)에 따라 OT message 순서를
바꾼다.

\[
  b^{owner}=0: (m_0,m_1)=(-r, x-r)
\]

\[
  b^{owner}=1: (m_0,m_1)=(x-r, -r)
\]

Choice owner는 자기 Boolean share \(b^{choice}\)를 OT choice bit로 사용한다. 그러면
네 가지 경우는 다음과 같이 정리된다.

\begin{table}[h]
\centering
\begin{tabular}{cccc}
\toprule
\(b^{owner}\) & \(b^{choice}\) & 실제 \(b\) & choice owner가 받는 값 \\
\midrule
0 & 0 & 0 & \(-r\) \\
0 & 1 & 1 & \(x-r\) \\
1 & 0 & 1 & \(x-r\) \\
1 & 1 & 0 & \(-r\) \\
\bottomrule
\end{tabular}
\end{table}

즉 선택된 메시지는 항상 다음과 같다.

\[
  (b^{owner}\oplus b^{choice})\cdot x-r.
\]

따라서 value owner의 share \(r\)와 choice owner의 share를 더하면 다음 식이 된다.

\[
  r + ((b^{owner}\oplus b^{choice})\cdot x-r)
  = (b^{owner}\oplus b^{choice})\cdot x \pmod p.
\]

\section*{두 번의 B2A-OT가 필요한 이유}

RsCpsi 이후 sender와 receiver는 각각 CPSI value share를 하나씩 가진다. 이를
\(w_i^S\), \(w_i^R\)라고 하면 교집합 위치에서

\[
  D_i^Y=w_i^S+w_i^R \pmod p
\]

이다. 따라서 전체 \(D_i^Y\)에 indicator \(b_i\)를 곱하려면 다음 두 값을 모두
share 형태로 만들어야 한다.

\[
  b_iw_i^S, \qquad b_iw_i^R.
\]

첫 번째 B2A에서는 sender가 value owner가 되고 receiver가 choice owner가 된다.
결과를 \((\alpha_i^S,\alpha_i^R)\)라고 하면 다음을 만족한다.

\[
  \alpha_i^S + \alpha_i^R = b_i \cdot w_i^S \pmod p.
\]

두 번째 B2A에서는 receiver가 value owner가 되고 sender가 choice owner가 된다.
결과를 \((\beta_i^S,\beta_i^R)\)라고 하면 다음을 만족한다.

\[
  \beta_i^S + \beta_i^R = b_i \cdot w_i^R \pmod p.
\]

각 참여자는 자신이 가진 두 share를 더한다.

\[
  a_i^S = \alpha_i^S + \beta_i^S \pmod p,
  \qquad
  a_i^R = \alpha_i^R + \beta_i^R \pmod p.
\]

따라서 최종적으로 다음이 성립한다.

\[
\begin{aligned}
  a_i^S+a_i^R
  &= b_i\cdot w_i^S + b_i\cdot w_i^R \\
  &= b_i\cdot (w_i^S+w_i^R) \\
  &= b_i\cdot D_i^Y \pmod p.
\end{aligned}
\]

\section*{Wrapper에서의 실행 순서}

Sender 쪽 wrapper는 먼저 RsCpsi sender를 실행한다. 그 후 sender가 가진 value share에
대해 value owner 역할로 B2A를 수행하고, receiver가 가진 value share에 대해서는 choice
owner 역할로 B2A를 수행한다.

\begin{lstlisting}[language=C++]
co_await send(identifiers, associatedData, share, chl);
co_await psiIpB2aValueOwner(share.mFlagBits, share.mValues,
    ownedValueShare, mConfig, mPrng, chl);
co_await psiIpB2aChoiceOwner(share.mFlagBits, share.mValues.rows(),
    share.mValues.cols(), receivedValueShare, mConfig, mPrng, chl);
psiIpAddPrimeShares(ownedValueShare, receivedValueShare,
    arithmeticShare, mConfig);
\end{lstlisting}

Receiver 쪽 wrapper는 같은 두 B2A를 반대 역할 순서로 수행한다.

\begin{lstlisting}[language=C++]
co_await receive(identifiers, share, chl);
co_await psiIpB2aChoiceOwner(share.mFlagBits, share.mValues.rows(),
    share.mValues.cols(), receivedValueShare, mConfig, mPrng, chl);
co_await psiIpB2aValueOwner(share.mFlagBits, share.mValues,
    ownedValueShare, mConfig, mPrng, chl);
psiIpAddPrimeShares(receivedValueShare, ownedValueShare,
    arithmeticShare, mConfig);
\end{lstlisting}

\section*{구현 세부사항}

\texttt{psiIpB2aValueOwner}는 \texttt{oc::MatrixView<u8>} 형태의 value share를 입력으로
받는다. matrix의 row는 CPSI output bin에 대응하고, column은 field element를 byte
단위로 저장한다. 기본 설정에서는 row마다 field element 하나가 있고, 각 element는
4 byte이다. 함수는 각 field element에 대해 OT 하나를 생성한다.

\texttt{psiIpB2aChoiceOwner}는 같은 row 수와 value byte length를 입력으로 받는다.
\texttt{mFlagBits}는 OT choice vector로 확장된다. row마다 field element가 여러 개라면
같은 Boolean share가 각 element에 반복되어 들어간다.

OT message는 \texttt{block} 타입으로 전달한다. 실제 prime element는 \texttt{u64}로
계산한 뒤 block의 하위 64 bit에 넣는다. 수신 후에는 다시 하위 64 bit를 읽고 \(p\)로
reduction하여 \(\mathbb{Z}_p\) 원소로 해석한다.

\texttt{psiIpAddPrimeShares}는 두 B2A 결과 matrix를 element 단위로 읽어서
\(\mathbb{Z}_p\) 덧셈을 수행한다. 이 함수의 output이 wrapper에서 반환하는
\texttt{arithmeticShare}이다.

\section*{검증}

검증용 테스트 \texttt{Cpsi\_PsiInnerproduct\_b2a\_test}를 추가하였다. 테스트에서는
일부 원소만 교집합이 되도록 입력을 만들고, B2A까지 끝난 후 양쪽 arithmetic share를
더해서 다음 값과 비교한다.

\[
  a_i^S + a_i^R =
  \begin{cases}
  D_i^Y & x_i \in Y,\\
  0 & x_i \notin Y.
  \end{cases}
\]

실행 명령은 다음과 같다.

\begin{lstlisting}[language=bash]
./out/build/osx/frontend/frontend -u Cpsi_PsiInnerproduct_b2a_test -n 64 -nt 1
\end{lstlisting}

테스트 결과는 통과하였다.

\end{document}
